% BEGIN VOL07-EXPANSION VII25
\section{Supplementary worked examples}
\label{sec:vii25-pedagogy}

\begin{example}[Trivial holonomy in Euclidean space]\label{ex:vii25-ped-01}
Every Euclidean loop has identity parallel transport for the standard Levi--Civita connection, so the holonomy group is trivial.  Flatness and a global parallel frame make path independence explicit.
\end{example}

\begin{example}[Sphere loop]\label{ex:vii25-ped-02}
Transport around a geodesic triangle on the unit sphere generally returns a vector rotated.  For a positively oriented geodesic triangle, the rotation angle equals its spherical area modulo \(2\pi\), a concrete Gauss--Bonnet manifestation of curvature as holonomy.
\end{example}

\begin{example}[Restricted versus global holonomy]\label{ex:vii25-ped-03}
A flat quotient can have locally trivial curvature while global topology affects transport for more general flat bundles.  This separates restricted holonomy, generated by contractible loops, from the full holonomy group, which can also record topology.
\end{example}

\section{Graded supplementary exercises}
The following exercises add a deliberate progression from concrete calculation to structural proof, hypothesis testing, application, and synthesis.

\subsection*{Standard computations and constructions}

\begin{exercise}[Group identity]\label{exr:vii25-ped-01}
Which loop gives the identity holonomy transformation?
\end{exercise}

\begin{hint}
Use the constant loop.
\end{hint}

\begin{solution}
The constant based loop transports every vector to itself.
\end{solution}

\begin{exercise}[Group inverse]\label{exr:vii25-ped-02}
Which loop gives the inverse of \(P_\gamma\)?
\end{exercise}

\begin{hint}
Reverse the loop.
\end{hint}

\begin{solution}
The reversed loop gives \(P_{\gamma^{-1}}=P_\gamma^{-1}\).
\end{solution}

\begin{exercise}[Euclidean holonomy]\label{exr:vii25-ped-03}
Compute \(\operatorname{Hol}_p(\mathbb R^n)\) for the standard metric.
\end{exercise}

\begin{hint}
Use path-independent Cartesian transport.
\end{hint}

\begin{solution}
It is the trivial group \(\{I\}\).
\end{solution}

\begin{exercise}[Orthogonal target]\label{exr:vii25-ped-04}
For Levi--Civita transport, in what group does holonomy lie?
\end{exercise}

\begin{hint}
Transport preserves inner products.
\end{hint}

\begin{solution}
\(\operatorname{Hol}_p\subseteq O(T_pM,g_p)\).
\end{solution}

\begin{exercise}[Sphere triangle]\label{exr:vii25-ped-05}
A geodesic triangle on the unit sphere has area \(A\).  What rotation angle is expected from Gauss--Bonnet holonomy?
\end{exercise}

\begin{hint}
Curvature is \(1\).
\end{hint}

\begin{solution}
The net rotation is \(A\) modulo \(2\pi\), with sign set by orientation conventions.
\end{solution}

\subsection*{Proofs}

\begin{exercise}[Group property]\label{exr:vii25-ped-06}
Prove holonomy transformations form a group.
\end{exercise}

\begin{hint}
Use constant loops, concatenation, and reversal.
\end{hint}

\begin{solution}
These operations provide identity, composition, and inverses, and transport is orthogonal.
\end{solution}

\begin{exercise}[Base-point conjugacy]\label{exr:vii25-ped-07}
Why do holonomy groups at two points in the same component differ only by conjugacy?
\end{exercise}

\begin{hint}
Transport loops along a connecting path.
\end{hint}

\begin{solution}
A path transport \(P_\sigma\) identifies tangent spaces and sends loop holonomy by \(H\mapsto P_\sigma H P_\sigma^{-1}\).
\end{solution}

\begin{exercise}[Restricted normality]\label{exr:vii25-ped-08}
Why is restricted holonomy normal in full holonomy?
\end{exercise}

\begin{hint}
Conjugate a contractible loop by an arbitrary loop.
\end{hint}

\begin{solution}
The conjugated loop remains contractible, so the subgroup generated by such loops is invariant under conjugation.
\end{solution}

\begin{exercise}[Parallel frame criterion]\label{exr:vii25-ped-09}
Prove a global parallel frame on a connected region forces trivial loop holonomy there.
\end{exercise}

\begin{hint}
Transport frame coefficients.
\end{hint}

\begin{solution}
Parallel transport leaves the constant coefficients in the parallel frame unchanged, hence every loop acts identically.
\end{solution}

\subsection*{Counterexamples and hypothesis tests}

\begin{exercise}[Flat implies trivial full holonomy]\label{exr:vii25-ped-10}
Test the claim: zero curvature always implies trivial full holonomy.
\end{exercise}

\begin{hint}
Distinguish local from global topology and general flat connections.
\end{hint}

\begin{solution}
False in general for flat connections on non-simply-connected spaces; flatness forces trivial restricted holonomy, while global monodromy may remain.
\end{solution}

\begin{exercise}[Holonomy is scalar]\label{exr:vii25-ped-11}
Test the claim: holonomy is a single angle in every dimension.
\end{exercise}

\begin{hint}
Consider \(O(n)\) for \(n>2\).
\end{hint}

\begin{solution}
False.  In higher dimensions holonomy transformations are general orthogonal maps, not one scalar angle.
\end{solution}

\begin{exercise}[Contractible loops irrelevant]\label{exr:vii25-ped-12}
Test the claim: only noncontractible loops can produce holonomy.
\end{exercise}

\begin{hint}
Use a curved sphere patch.
\end{hint}

\begin{solution}
False.  Small contractible loops already detect local curvature.
\end{solution}

\subsection*{Applications and investigations}

\begin{exercise}[Curvature sensor]\label{exr:vii25-ped-13}
How can a small-loop transport experiment serve as a curvature sensor?
\end{exercise}

\begin{hint}
Compare the returned vector with the initial one.
\end{hint}

\begin{solution}
The leading deviation from identity is controlled by curvature integrated over the enclosed small area.
\end{solution}

\begin{exercise}[Frame consistency]\label{exr:vii25-ped-14}
Why does nontrivial holonomy obstruct globally path-independent frame propagation?
\end{exercise}

\begin{hint}
Transport a frame around a loop.
\end{hint}

\begin{solution}
Returning with a different frame means two paths to the same point cannot define a single consistent propagated orientation.
\end{solution}

\subsection*{Challenge problems}

\begin{exercise}[Abelian surface holonomy]\label{exr:vii25-ped-15}
On an oriented two-dimensional Riemannian surface, why can local holonomy be described by a rotation angle?
\end{exercise}

\begin{hint}
The oriented orthogonal group in two dimensions is \(SO(2)\).
\end{hint}

\begin{solution}
Orientation-preserving Levi--Civita transport acts by elements of \(SO(2)\), each determined by one angle.
\end{solution}

\begin{exercise}[Flat simply connected]\label{exr:vii25-ped-16}
Explain why a simply connected flat Riemannian manifold has trivial restricted and, locally under completeness assumptions, path-independent Levi--Civita transport.
\end{exercise}

\begin{hint}
All loops are contractible and curvature vanishes.
\end{hint}

\begin{solution}
Restricted holonomy is trivial by flatness; simple connectivity removes additional loop classes, so loop transport is trivial and transport depends only on endpoints.
\end{solution}

% END VOL07-EXPANSION VII25
